\documentclass[11pt]{article}

\usepackage[a4paper,margin=1in]{geometry}
\usepackage{amsmath,amssymb,amsthm,mathtools}
\usepackage{hyperref}

\newtheorem{theorem}{Theorem}
\newtheorem{proposition}{Proposition}
\newtheorem{remark}{Remark}
\newtheorem{definition}{Definition}

\DeclareMathOperator{\BL}{BL}
\DeclareMathOperator{\tr}{tr}
\DeclareMathOperator{\diag}{diag}

\title{Brascamp--Lieb Inequalities, Young's Inequality, and Optimal Transport}
\author{}
\date{}

\begin{document}

\maketitle

\section{The Brascamp--Lieb inequality}

The Brascamp--Lieb inequality is a common generalisation of several classical
integral inequalities, including H\"older's inequality, the Loomis--Whitney
inequality, and Young's convolution inequality.

\begin{definition}[Brascamp--Lieb datum]
A Brascamp--Lieb datum consists of surjective linear maps
\[
        B_j : \mathbb R^n \longrightarrow \mathbb R^{n_j},
        \qquad j=1,\dots,m,
\]
together with non-negative exponents
\[
        p_j \geq 0.
\]
For non-negative integrable functions
\[
        f_j : \mathbb R^{n_j} \to [0,\infty),
\]
the Brascamp--Lieb functional is
\[
        \mathcal B_{\mathbf B,\mathbf p}(f_1,\dots,f_m)
        :=
        \frac{
        \displaystyle
        \int_{\mathbb R^n}
        \prod_{j=1}^m f_j(B_jx)^{p_j}\,dx
        }{
        \displaystyle
        \prod_{j=1}^m
        \left(\int_{\mathbb R^{n_j}} f_j\right)^{p_j}
        }.
\]
The optimal Brascamp--Lieb constant is
\[
        \BL(\mathbf B,\mathbf p)
        :=
        \sup_{f_j\geq 0}
        \mathcal B_{\mathbf B,\mathbf p}(f_1,\dots,f_m).
\]
Thus the Brascamp--Lieb inequality is
\[
        \int_{\mathbb R^n}
        \prod_{j=1}^m f_j(B_jx)^{p_j}\,dx
        \leq
        \BL(\mathbf B,\mathbf p)
        \prod_{j=1}^m
        \left(\int_{\mathbb R^{n_j}} f_j\right)^{p_j}.
\]
\end{definition}

The finiteness of the constant is governed by two geometric conditions. The
first is the scaling condition
\[
        n = \sum_{j=1}^m p_j n_j.
\]
This is necessary by dilating all the functions simultaneously. The second is
the dimension condition
\[
        \dim V
        \leq
        \sum_{j=1}^m p_j \dim B_jV
        \qquad
        \text{for every subspace } V\leq \mathbb R^n.
\]
This prevents mass from concentrating along a subspace on which the right-hand
side has insufficient dimension to control the left-hand side.

\begin{theorem}[Finiteness criterion]
For a non-degenerate Brascamp--Lieb datum, the constant
\(\BL(\mathbf B,\mathbf p)\) is finite if and only if
\[
        n = \sum_{j=1}^m p_j n_j
\]
and
\[
        \dim V
        \leq
        \sum_{j=1}^m p_j \dim B_jV
        \qquad
        \text{for every subspace } V\leq \mathbb R^n.
\]
\end{theorem}

\section{Lieb's Gaussian formula}

A central theorem of Lieb says that the optimal Brascamp--Lieb constant is
already detected by Gaussian functions.

For a positive definite matrix \(A_j\) on \(\mathbb R^{n_j}\), define the
centred Gaussian
\[
        G_{A_j}(y)
        :=
        e^{-\pi \langle A_jy,y\rangle}.
\]
Then
\[
        \int_{\mathbb R^{n_j}} G_{A_j}(y)\,dy
        =
        (\det A_j)^{-1/2}.
\]
On the other hand,
\[
        \prod_{j=1}^m G_{A_j}(B_jx)^{p_j}
        =
        \exp\left(
        -\pi
        \left\langle
        \left(\sum_{j=1}^m p_j B_j^*A_jB_j\right)x,x
        \right\rangle
        \right).
\]
Therefore
\[
        \int_{\mathbb R^n}
        \prod_{j=1}^m G_{A_j}(B_jx)^{p_j}\,dx
        =
        \det\left(
        \sum_{j=1}^m p_j B_j^*A_jB_j
        \right)^{-1/2}.
\]
The value of the Brascamp--Lieb functional on these Gaussians is
\[
        \left[
        \frac{
        \prod_{j=1}^m (\det A_j)^{p_j}
        }{
        \det\left(
        \sum_{j=1}^m p_j B_j^*A_jB_j
        \right)
        }
        \right]^{1/2}.
\]

\begin{theorem}[Lieb's Gaussian theorem]
The sharp Brascamp--Lieb constant is given by the finite-dimensional Gaussian
optimisation problem
\[
        \BL(\mathbf B,\mathbf p)
        =
        \sup_{A_j>0}
        \left[
        \frac{
        \prod_{j=1}^m (\det A_j)^{p_j}
        }{
        \det\left(
        \sum_{j=1}^m p_j B_j^*A_jB_j
        \right)
        }
        \right]^{1/2}.
\]
In particular, whenever the supremum is attained by positive definite matrices
\(A_j\), the corresponding centred Gaussians
\[
        G_{A_j}(y)=e^{-\pi\langle A_jy,y\rangle}
\]
are extremisers for the Brascamp--Lieb inequality.
\end{theorem}

In degenerate cases the supremum need not be attained, but Lieb's theorem still
says that Gaussian extremising sequences determine the optimal constant.

\section{Young's inequality}

Young's convolution inequality is one of the model examples contained in the
Brascamp--Lieb framework.

\begin{definition}[Young's convolution inequality]
Let \(1\leq p,q,r\leq \infty\) satisfy
\[
        \frac1p+\frac1q = 1+\frac1r.
\]
Then
\[
        \|F*G\|_{L^r(\mathbb R^n)}
        \leq
        \|F\|_{L^p(\mathbb R^n)}
        \|G\|_{L^q(\mathbb R^n)}.
\]
The sharp form, due to Beckner and Brascamp--Lieb, is
\[
        \|F*G\|_{L^r(\mathbb R^n)}
        \leq
        \left(\frac{A_pA_q}{A_r}\right)^n
        \|F\|_{L^p(\mathbb R^n)}
        \|G\|_{L^q(\mathbb R^n)},
\]
where
\[
        A_s
        :=
        s^{1/(2s)}(s')^{-1/(2s')},
        \qquad
        \frac1s+\frac1{s'}=1,
\]
with the usual endpoint convention \(A_1=A_\infty=1\).
\end{definition}

To see Young's inequality as a Brascamp--Lieb inequality, use duality. For
\(1<r<\infty\),
\[
        \|F*G\|_r
        =
        \sup_{\|H\|_{r'}=1}
        \int_{\mathbb R^n} (F*G)(z)H(z)\,dz.
\]
Writing out the convolution gives the trilinear form
\[
        \int_{\mathbb R^n}\int_{\mathbb R^n}
        F(x)G(y)H(x+y)\,dx\,dy.
\]
This is a Brascamp--Lieb expression on the ambient space
\(\mathbb R^n\times \mathbb R^n\), with linear maps
\[
        B_1(x,y)=x,
        \qquad
        B_2(x,y)=y,
        \qquad
        B_3(x,y)=x+y.
\]
Set
\[
        c_1=\frac1p,
        \qquad
        c_2=\frac1q,
        \qquad
        c_3=\frac1{r'}.
\]
If
\[
        f_1=F^p,\qquad f_2=G^q,\qquad f_3=H^{r'},
\]
then
\[
        f_1(B_1(x,y))^{c_1}
        f_2(B_2(x,y))^{c_2}
        f_3(B_3(x,y))^{c_3}
        =
        F(x)G(y)H(x+y).
\]
The Brascamp--Lieb inequality therefore gives
\[
        \int_{\mathbb R^n}\int_{\mathbb R^n}
        F(x)G(y)H(x+y)\,dx\,dy
        \leq
        C
        \|F\|_p\|G\|_q\|H\|_{r'}.
\]
The scaling condition for this Brascamp--Lieb datum is
\[
        2n
        =
        n(c_1+c_2+c_3),
\]
which is equivalent to
\[
        \frac1p+\frac1q+\frac1{r'}=2.
\]
Since \(1/r'=1-1/r\), this is exactly
\[
        \frac1p+\frac1q=1+\frac1r.
\]
Thus Young's inequality is precisely the Brascamp--Lieb inequality associated
with the three maps
\[
        (x,y)\mapsto x,\qquad
        (x,y)\mapsto y,\qquad
        (x,y)\mapsto x+y.
\]

The sharp Young constant is the Gaussian Brascamp--Lieb constant for this datum:
\[
        C
        =
        \left(\frac{A_pA_q}{A_r}\right)^n.
\]
Equivalently,
\[
        C
        =
        \left(A_pA_qA_{r'}\right)^n,
\]
since \(A_{r'}=A_r^{-1}\).

\section{Optimal transport and the geometric Brascamp--Lieb inequality}

The connection with optimal transport is clearest after putting the datum in
geometric position.

\begin{definition}[Geometric Brascamp--Lieb datum]
A Brascamp--Lieb datum is in geometric position if the maps
\[
        U_j:\mathbb R^n\to \mathbb R^{n_j}
\]
satisfy
\[
        U_jU_j^*=I_{\mathbb R^{n_j}}
\]
and
\[
        \sum_{j=1}^m p_jU_j^*U_j=I_{\mathbb R^n}.
\]
\end{definition}

In this case the Brascamp--Lieb inequality has constant \(1\):
\[
        \int_{\mathbb R^n}
        \prod_{j=1}^m f_j(U_jx)^{p_j}\,dx
        \leq
        \prod_{j=1}^m
        \left(\int_{\mathbb R^{n_j}}f_j\right)^{p_j}.
\]

The standard Gaussian already gives equality. Indeed, with
\[
        \gamma_k(x)=e^{-\pi |x|^2},
\]
we have \(\int_{\mathbb R^k}\gamma_k=1\), and
\[
        \prod_{j=1}^m \gamma_{n_j}(U_jx)^{p_j}
        =
        \exp\left(
        -\pi\sum_{j=1}^m p_j|U_jx|^2
        \right).
\]
Using
\[
        \sum_{j=1}^m p_jU_j^*U_j=I_{\mathbb R^n},
\]
we get
\[
        \sum_{j=1}^m p_j|U_jx|^2=|x|^2.
\]
Hence
\[
        \int_{\mathbb R^n}
        \prod_{j=1}^m \gamma_{n_j}(U_jx)^{p_j}\,dx
        =
        \int_{\mathbb R^n}e^{-\pi |x|^2}\,dx
        =
        1.
\]

Now let \(f_j\) be smooth positive probability densities on
\(\mathbb R^{n_j}\). Let \(T_j=\nabla\phi_j\) be the Brenier map transporting
\(f_j(y)\,dy\) optimally, for quadratic cost, to the standard Gaussian measure
\(\gamma_{n_j}(z)\,dz\). The Monge--Ampere equation gives
\[
        f_j(y)
        =
        \gamma_{n_j}(T_j(y))\det DT_j(y).
\]
Define
\[
        S(x)
        :=
        \sum_{j=1}^m p_jU_j^*T_j(U_jx).
\]
Since each \(T_j\) is the gradient of a convex function, \(S\) is also the
gradient of a convex function:
\[
        S=\nabla\left(\sum_{j=1}^m p_j\phi_j(U_jx)\right).
\]
Its Jacobian is
\[
        DS(x)
        =
        \sum_{j=1}^m
        p_jU_j^*DT_j(U_jx)U_j.
\]

Two inequalities are the heart of Barthe's transport proof. First, by the
geometric condition and Cauchy's inequality,
\[
        |S(x)|^2
        \leq
        \sum_{j=1}^m p_j|T_j(U_jx)|^2.
\]
Second, by the determinant inequality associated with geometric
Brascamp--Lieb data,
\[
        \det DS(x)
        \geq
        \prod_{j=1}^m
        \det DT_j(U_jx)^{p_j}.
\]
Combining these with the Monge--Ampere identities yields
\[
\begin{aligned}
        \prod_{j=1}^m f_j(U_jx)^{p_j}
        &=
        \prod_{j=1}^m
        \gamma_{n_j}(T_j(U_jx))^{p_j}
        \prod_{j=1}^m
        \det DT_j(U_jx)^{p_j}    \\
        &=
        \exp\left(
        -\pi\sum_{j=1}^m p_j|T_j(U_jx)|^2
        \right)
        \prod_{j=1}^m
        \det DT_j(U_jx)^{p_j}    \\
        &\leq
        e^{-\pi |S(x)|^2}\det DS(x)                 \\
        &=
        \gamma_n(S(x))\det DS(x).
\end{aligned}
\]
Since \(S\) is a gradient map, the area formula gives
\[
        \int_{\mathbb R^n}\gamma_n(S(x))\det DS(x)\,dx
        \leq
        \int_{\mathbb R^n}\gamma_n(z)\,dz
        =
        1.
\]
Therefore
\[
        \int_{\mathbb R^n}
        \prod_{j=1}^m f_j(U_jx)^{p_j}\,dx
        \leq 1.
\]
By homogeneity this gives the full geometric Brascamp--Lieb inequality.

\section{From the general case to the geometric case}

Suppose the Gaussian optimisation problem is attained by positive definite
matrices \(A_j\), and set
\[
        M
        :=
        \sum_{j=1}^m p_jB_j^*A_jB_j.
\]
Define
\[
        U_j
        :=
        A_j^{1/2}B_jM^{-1/2}.
\]
Then
\[
        \sum_{j=1}^m p_jU_j^*U_j
        =
        I_{\mathbb R^n}.
\]
Moreover, the Euler--Lagrange equations for the Gaussian determinant problem
give
\[
        A_j^{-1}=B_jM^{-1}B_j^*.
\]
Consequently,
\[
        U_jU_j^*
        =
        A_j^{1/2}B_jM^{-1}B_j^*A_j^{1/2}
        =
        I_{\mathbb R^{n_j}}.
\]
Thus the Gaussian extremiser puts the datum into geometric position. In that
position, the standard Gaussian extremises the inequality with constant \(1\).
Undoing the linear changes of variables gives the extremising Gaussians
\[
        G_{A_j}(y)=e^{-\pi\langle A_jy,y\rangle}
\]
for the original datum.

\section{Why Gaussians are maximisers}

There are three complementary ways to understand the role of Gaussians.

First, Lieb's theorem gives the exact statement:
\[
        \BL(\mathbf B,\mathbf p)
        =
        \sup_{\text{centred Gaussians}}
        \mathcal B_{\mathbf B,\mathbf p}.
\]
Thus no non-Gaussian family of functions can produce a larger value than the
best Gaussian family.

Second, Gaussians are algebraically stable under the operations appearing in
Brascamp--Lieb inequalities. If
\[
        f_j(y)=e^{-\pi\langle A_jy,y\rangle},
\]
then
\[
        \prod_{j=1}^m f_j(B_jx)^{p_j}
\]
is again a Gaussian in \(x\), because the exponent is a quadratic form:
\[
        \sum_{j=1}^m p_j\langle A_jB_jx,B_jx\rangle
        =
        \left\langle
        \left(\sum_{j=1}^m p_jB_j^*A_jB_j\right)x,x
        \right\rangle.
\]
Therefore the infinite-dimensional variational problem over all functions
collapses, for Gaussians, to a finite-dimensional determinant problem.

Third, the optimal-transport proof explains why Gaussians are the equality
case. In geometric position, arbitrary densities \(f_j\) are transported to
Gaussians by Brenier maps \(T_j\). The proof compares the original integral to
\[
        \int_{\mathbb R^n}\gamma_n(S(x))\det DS(x)\,dx.
\]
The comparison uses two inequalities: a quadratic Cauchy inequality and a
determinant inequality. When the \(f_j\) are Gaussians, the Brenier maps are
linear; in geometric position they are simply the identity maps. Then both
inequalities become equalities. For genuinely non-Gaussian inputs, the Brenier
maps are nonlinear, and the determinant and quadratic steps create loss, except
in cases forced by the symmetries or factorisations of the datum.

For Young's inequality this recovers the classical sharp statement: in the
strict range \(1<p,q,r<\infty\), the extremisers for the sharp Young inequality
are Gaussian functions, up to the natural symmetries of translation, dilation,
multiplication by constants, and common linear changes of variables.

\section{Entropy viewpoint}

There is also an entropy formulation. Let \(X\) be a random vector in
\(\mathbb R^n\) with density \(\rho\), and let \(B_jX\) have density \(\rho_j\).
The Brascamp--Lieb inequality is dual to the entropy inequality
\[
        h(X)
        \leq
        \sum_{j=1}^m p_j h(B_jX)
        +
        \log \BL(\mathbf B,\mathbf p),
\]
where
\[
        h(X)
        =
        -\int_{\mathbb R^n}\rho(x)\log\rho(x)\,dx
\]
is differential entropy. This is another place where Gaussians naturally enter:
among all random variables with fixed covariance, Gaussians maximise entropy.
Thus the Gaussian extremisers are consistent with both the determinant formula
and the transport/entropy interpretation.

\begin{thebibliography}{9}

\bibitem{BrascampLieb}
H. J. Brascamp and E. H. Lieb,
\emph{Best constants in Young's inequality, its converse, and its generalization to more than three functions},
Advances in Mathematics, 20 (1976), 151--173.

\bibitem{Beckner}
W. Beckner,
\emph{Inequalities in Fourier analysis},
Annals of Mathematics, 102 (1975), 159--182.

\bibitem{Lieb}
E. H. Lieb,
\emph{Gaussian kernels have only Gaussian maximizers},
Inventiones Mathematicae, 102 (1990), 179--208.

\bibitem{Barthe}
F. Barthe,
\emph{On a reverse form of the Brascamp--Lieb inequality},
Inventiones Mathematicae, 134 (1998), 335--361.

\bibitem{BCCT}
J. Bennett, A. Carbery, M. Christ and T. Tao,
\emph{The Brascamp--Lieb inequalities: finiteness, structure and extremals},
Geometric and Functional Analysis, 17 (2008), 1343--1415.

\end{thebibliography}

\end{document}